\section{引言}\label{sec:intro}

一个借助 Python 解释器解数学题、上网检索、或在数月间维护用户记忆的 LLM 智能体（agent），要经历多轮（turn）行动：写代码、读工具输出、决定存储或删除什么，直到最后才收到一个数字——对或错。强化学习（RL）随后必须判断这些轮中哪些该得到功劳。2025 年的主流算法 GRPO \citep{shao2024deepseekmath} 每次给出的答案都一样：所有轮，一视同仁。一次决定性的工具调用、一次多余的调用和一个格式 token 得到相同的优势（advantage），正确的归因只能在大量样本上以统计方式浮现。这就是\emph{信用分配}（credit assignment）问题，RL 中最古老的问题之一 \citep{sutton1988learning,pignatelli2024survey}，如今它以一种更新、更尖锐的形式回来了。

\paragraph{为什么是现在}
三方面的进展让这个问题变得紧迫。第一，可验证奖励强化学习（RLVR）在一年之内从单回复推理扩展到了多轮智能体：数学、搜索、软件工程和深度研究智能体如今都用结果奖励端到端训练 \citep{zhang2025landscape}，horizon（时域）长达数百次工具调用，统一分配的信用在这种设置下越来越嘈杂。第二，主导 2025 年的稳定性问题（熵坍缩、工具引起的漂移、零方差组（zero-variance group））如今大多已有成熟配方解决，剩下的问题就是归因本身；2025--2026 年出现的轮级信用方法远不止十几种，大多只与 GRPO 比较。第三，记忆智能体暴露了标准组内基线无法处理的问题。Memory-R2 \citep{yan2026memory} 指出，一旦各条 rollout 写入、更新或删除了不同的记忆，它们面对的就不再是同一个环境，比较它们的最终奖励“从根本上是不公平的”\unverified。它的 LoGo-GRPO 从共享的记忆快照出发增加局部重 rollout（rerollout）。我们把它作为贯穿全文的例子，因为它同时展示了基于分支的信用方法的潜力，以及当前证据的主要缺口：报告的增益没有等算力（compute-matched）基线。

\paragraph{范围与方法}
我们关注 2025--2026 年用 RL 训练的 LLM 智能体的信用分配工作，以及这些方法所沿用的经典 RL 思想（TD 学习、RUDDER、事后（hindsight）与反事实信用）。资料主要是截至 2026 年 11 月收集的 arXiv 预印本、项目主页和技术博客；不少 2026 年的论文我们只能看到摘要、HTML 节选或第三方总结。凡是无法对照原始来源核实的说法（数字、arXiv 编号、发表会议、作者列表），我们都用剑号\unverified{}标出，原文中的保留措辞也一律保留。读者应把带剑号的内容当作待查线索，而非既定事实。除非另有说明，“增益”指在相同轨迹（trajectory）数下相对 GRPO 的提升——这是大多数论文的报告方式，而我们认为这往往是错误的比较。本综述还用到我们自己的一项对照研究，即 ToolCredit 研究，其结果大多为阴性，我们如实报告。

\paragraph{贡献}
\begin{itemize}
  \item 一个统一视角：把 GRPO 家族算法看作\emph{统一}信用，把智能体 RL 的各种配方看作不改变归因的稳定性修补（第~\ref{sec:background}--\ref{sec:agentic}~节）。
  \item 将智能体信用方法分为六大类的分类体系，按其对环境的要求（重置、重复出现的状态、验证器、标准答案）和额外算力的去向逐一刻画，并给出按环境性质选择方法的指南（第~\ref{sec:taxonomy}--\ref{sec:comparison}~节）。
  \item 对 Memory-R2 与 LoGo-GRPO 的细致批判性解读，区分哪些是证据已经表明的，哪些仍需要等算力基线来确认（第~\ref{sec:memory}~节）。
  \item 来自一项预注册、等算力的数学工具调用轮级信用研究的教训：非守恒信用充当了奖励并被钻空子；守恒约束消除了这一病态，但在 $\leq$4 轮时没有可检测的增益；DAPO 式过滤在每条 rollout 上没有带来增益（第~\ref{sec:toolcredit}~节）。
  \item 一份研究议程，包含按优先级排序并附预算的想法和一个十周计划；首要建议是构建一个 horizon 可控、真实信用可计算的基准（第~\ref{sec:agenda}~节）。
\end{itemize}

\paragraph{全文结构}
第~\ref{sec:background}~节回顾 RLHF、RLVR 和 GRPO 家族，说明它们的改进都让信用保持统一。第~\ref{sec:agentic}~节介绍多轮智能体 RL、其任务类别，以及为什么能否重置到中间状态决定了哪些信用方法可行。第~\ref{sec:problem}~节阐述信用分配问题、它的经典工具，以及只有结果奖励时的统计代价。第~\ref{sec:taxonomy}~节把现有方法归为六大类，并回顾它们的失效模式与审查。第~\ref{sec:memory}~节转向记忆智能体和 Memory-R2。第~\ref{sec:comparison}~节比较各类方法并给出选择指南。第~\ref{sec:toolcredit}~节报告 ToolCredit 研究，第~\ref{sec:agenda}~节把以上内容转化为研究方向和计划。附录给出阅读清单和术语表。
